\section{A double cover joined to the zero section}\label{sec:cubic}
For a double wall the first reducible support occurs already in degree three. A branch point supplies a joining flag, and different flags may occur together. This leads to a relation between the signed counts of smooth double and cubic covers.
\subsection{Smoothing a double cover joined to the identity}\label{p10-sec:nodes}

Let $D_t$ be a smooth path of real Cauchy--Riemann operators with fixed
symbol on a complex rank-two bundle $N$ of degree $2g-2$ over $C$.
Suppose $D=D_0$ is invertible, of determinant sign $\epsilon$.
Let $f:Y\to C$ be a connected smooth double cover with $b>0$ simple
ramification points over distinct values. Suppose that $K=\ker f^*D$
is a real line, that its generator $\xi$ is odd under the deck
involution, and that $(f,\xi):Y\to N$ is an embedding. Put
$P_f=\operatorname{coker}f^*D$. Write $s$ for the $b$ complex branch
parameters and the increasing real path parameter together. We assume
the double crossing is transverse, so that the derivative of the
pullback operator on its kernel induces an isomorphism
\[
  W_f:T_{(f,0)}(\mathcal H_2\times\RR)
          \longrightarrow \operatorname{Hom}_{\RR}(K,P_f).
\]
Here $\mathcal H_2$ denotes the fixed-target double Hurwitz chart.
Write $\eta_f$ for the sign of $W_f$ with the absolute Fredholm
orientation, obtained by deformation to a complex operator.

Choose a nonempty subset $I$ of the ramification points. Join each
$r_i\in I$ to $x_i=f(r_i)$ on one copy of $C$, and denote the resulting
finite cubic map by $h_I$. Its arithmetic genus is
$H=g(Y)+g+j-1$, where $j=|I|$. Near $r_i$ take coordinates
$w=z^2$ centered at $x_i,r_i$, and a local frame holomorphic for the
complex-linear part of $D$. Write $D=\bar\partial+B(w)\,d\bar w$ with
$B(w)$ antilinear, and put $v_i=\partial_z\xi(r_i)$. Since $df(r_i)=0$,
$d\xi(r_i)$ is complex linear, and embeddedness gives $v_i\ne0$.
The condition
\begin{equation}\label{p10-eq:node-jet}
                 B(x_i)v_i\ne 0\quad\text{in }N_{x_i}/\CC v_i
\end{equation}
is independent of the coordinates, complex frame and real generator.
Indeed the change of frame adds only a complex-linear term to the
zeroth-order coefficient, and $\xi(r_i)=0$.

\begin{theorem}\label{p10-thm:nodes}
With the double crossing and joining set $I$ as above, assume \eqref{p10-eq:node-jet} at every point
of $I$. The pullback kernel at $h_I$ is isolated in the full fixed-target
chart of genus $H$ times the path. Its local contribution to the change
of the absolutely oriented normal Euler number is
\[
                         (-1)^j\epsilon\eta_f .
\]
If \eqref{p10-eq:node-jet} holds at every ramification point, summing over
all choices of $j$ joins gives $\binom bj(-1)^j\epsilon\eta_f$.
Condition \eqref{p10-eq:node-jet} holds at all such crossings on a generic
smooth path of base-invertible fixed-symbol operators with rigid endpoints.
\end{theorem}

The attaching parameter moves the common value along the tangent of
the double graph. The smoothing parameter contributes in the
quotient by that tangent. Its leading term is conjugate-linear after
a radial change of variable, which accounts for the sign at a node.

The full local calculation is given in Appendix~\ref{app:double-equation}.
Its normalized equation has one complex tangent column from the
moving branch and one complex antilinear column from the smoothing
at each selected point. The projected coefficient in
\eqref{p10-eq:node-jet} makes both columns regular after the radial
change of variable. The base determinant and the absolute crossing
coorientation supply the two remaining signs. Unlike the sextic
calculation below, the double has only an odd kernel line; imposing
its normalization leaves no even height common to the nodes.

To pass from this local degree to the Gromov--Witten change, one must
also allow contracted curves at the joins and at smooth points of
the two branches. The displacement comparison of
Theorem~\ref{p10-thm:descendants} replaces the joined map by regular
embedded graphs of genus \(H\). In source genus \(k\) their degree-one
contribution is \([v^{2k-2}]x^{H-1}\) times the signed degree above.
Proposition~\ref{p10-prop:compact} gives the finiteness needed to sum
these contributions along a generic path.

\subsection{The chamber formula and the balanced product}\label{p10-sec:chamberproof}
\begin{proof}[Proof of Theorem~\ref{p10-thm:chambers}]
By Proposition~\ref{p10-prop:compact} and its double-cover counterpart
there are finitely many crossings on the path. At a smooth double
wall with $2l$ branch values the degree-two change is $\eta_f x^{2c+l}$
(Theorem~\ref{p9-thm:wall}). At a smooth cubic wall with the same
branch count the degree-three change is $2\eta_f x^{3c+l}$ by
Proposition~\ref{p10-prop:smoothjump}. The separation in
Proposition~\ref{p10-prop:incidence} applies to covers of every genus.

Consider the degree-three virtual cylinder at a double crossing.
A nonzero section must restrict nontrivially to a nonconstant
normalization, since matching determines its values on every
connected contracted block. There is no cubic normalization with
kernel at this parameter. There must therefore be a degree-two
normalization carrying the odd kernel and a copy of $C$ at zero
height. Both graphs are embedded. Their common image points are the
$2l$ ramification points, each with one preimage on each graph.
There is no self-node and a contracted block has at most two
nonconstant attaching branches. Thus at least one of these $2l$
points is used to connect the graphs, either directly or by a
contracted bridge. The set $I$ of such points determines the compact
locus of Theorem~\ref{p10-thm:descendants}; every other contracted
component belongs to its allowed decorations. When $l=0$ there
is no connected support of this type.

Put $j=|I|$. The arithmetic genus of the bare join is
$(1+2c+l)+(1+c)+j-1=1+3c+l+j$. Theorem~\ref{p10-thm:descendants}
therefore gives the total change at this double crossing as
\[
  \sum_{j=1}^{2l}\binom{2l}{j}(-1)^j\epsilon\eta_f x^{3c+l+j}
   =\epsilon\eta_f x^{3c+l}\big((1-x)^{2l}-1\big).
\]
For each fixed source genus the supported Euler change is additive
on the finite collection of disjoint compact loci just classified.
Outside them the ordinary oriented cokernel extends over the
parameter. Summing coefficientwise over all arithmetic genera proves
\eqref{p10-eq:chambers}. Substitution $x\mapsto (2\sin(v/2))^2$ is
injective on polynomials, since the substituted series begins with
$v^2$. The first equation determines $W_2$ from the endpoints and
the second then determines $W_3$.
\end{proof}

Let $H$ be Pardon's rational algebra of complex cycle families
\cite{pardon}, and $F_h$ the balanced antidiagonal local curve series
of genus $h$ in this algebra \cite{WueProduct}. In the comparison
algebra $S$ of almost complex cycle families, let $P_d(D)$ denote
the class of the isolated cycle $d[C]$ for the homogeneous normal
equation $D$. The variable $q$ records the degree over $C$. A
\emph{balanced lift through degree three} is a product of
$F_h(q^d)$ with integer exponents, finite at each $d\le3$, whose
image under the comparison $H\to S$ agrees with
$1+\sum_{d=1}^3P_d(D)q^d$ modulo $q^4$.

\begin{theorem}\label{p10-thm:product}
Each endpoint in Theorem~\ref{p10-thm:chambers} has a unique balanced lift
$P^H_{\le3}(D)$ through degree three. Set
$V(x)=2W_3(x)+\epsilon(W_2(x(1-x)^2)-W_2(x))$. Then
\[
 \frac{P^H_{\le3}(D_+)}{P^H_{\le3}(D_-)}
 =\prod_{l\ge0}F_{1+2c+l}(q^2)^{W_{2,l}}\,
  \prod_{l\ge0}F_{1+3c+l}(q^3)^{[x^l]V(x)}
       \pmod{q^4}.
\]
\end{theorem}
\begin{proof}
We apply the finite-degree factorization of \cite[Corollary~3.2]{WueUnramified}
using only the stated endpoint rigidity. Smooth the intersections of a
horizontal torus with $g-1$ parallel vertical tori in a product of
two symplectic two-tori, and include the product in a six-torus. The
resulting genus-$g$ surface represents a primitive class $A$,
as pairing in the first two factors with a vertical torus gives
one. Its symplectic normal bundle has degree $2g-2$. Identify the
normal germ with the given one along the surface and extend its graph
almost complex structure to a global tame structure, by interpolation
to a compatible extension in a sufficiently small tube
\cite[proof of Proposition~3.1]{WueUnramified}.

Through energy $3\omega(A)$ the only cycles contained in the
closed tube are multiples of the zero section. Indeed a component
in that tube projects holomorphically in positive degree; a component
of nonzero height would give a nonzero kernel on its smooth
normalization. The degree agrees with the positive energy label since
the tube retracts to $C$. The cycles in the interior give a compact
open and closed insertion through this cutoff. Compactness of cycles
and convergence of supports preserve separation from the tube
boundary under a small perturbation of the ambient structure
\cite{diw}. Transport to a generic perturbation and finite-energy
cluster decomposition express this insertion as a product of the
images of the balanced factors with integral exponents
\cite[Corollary~3.2]{WueUnramified}. Their coefficients generate the integral
cycle algebra by \cite[Theorem~1.1]{WueGenerators}. The construction at that energy involves
only labels through three. Lift each factor to $H$. If two products
first differed at label $d$, the difference at that label would be
a finite linear combination of $[q]F_h(q)$ whose image vanishes
in $S$. Independence of these images under fixed-genus detection
\cite[Theorem~5.1]{WueProduct} gives
uniqueness.

It remains to identify the exponents of these geometric products.
At each endpoint, fixed-target comparison in the compact realization
\cite[Theorem~1.2]{zinger} identifies the connected local series
in degrees $1,2,3$ with the integrals defining $L_d$. The series
of an elementary factor in its first label is $x^{h-1}$
\cite[Theorem~2]{pandhodge}. In degree one the endpoint factor is
$F_g(q)^{\epsilon}$ at both ends. At labels two and three the only
proper divisor is one. Its contribution to the logarithm of the
ratio cancels, and an exponent difference in genus $h$ at either
new label therefore contributes exactly its multiple of $x^{h-1}$.
The finite Laurent monomials in $x$ are linearly independent after
sine substitution. Equation~\eqref{p10-eq:chambers} determines all
the differences asserted above.
\end{proof}

Write $n_{d,h}(D)$ for the exponents in this lift. The original
fixed-target series in degree $d$ vanishes below genus $1+dc$ by
Riemann--Hurwitz. For $d=2,3$ the only proper-divisor contribution
comes from the genus-$g$ balanced reference in label one, whose
degree-$d$ term satisfies the same bound. The remaining monomials
$x^{h-1}$ have distinct leading orders, so $n_{d,h}(D)=0$ for
$h<1+dc$ \cite[Proposition~3.1 and Corollary~3.2]{WueUnramified}. Set
$N_d(D;x)=\sum_{l\ge0} n_{d,1+dc+l}(D)x^l$.
\begin{corollary}
The polynomial $N_3(D;x)-N_2(D;x+x^3)+N_2(D;x)$ modulo two is
constant between rigid endpoints joined in $\mathcal U$.
\end{corollary}
\begin{proof}
The two exponent changes are $\Delta N_2=W_2$ and $\Delta N_3=V$.
Modulo two one has $x(1-x)^2=x+x^3$, $2W_3=0$ and $\epsilon=1$.
\end{proof}
